\documentclass[11pt,a4paper]{article}

% ---------- fontes e língua ----------
\usepackage[T1]{fontenc}
\usepackage[utf8]{inputenc}
\usepackage{lmodern}
\usepackage{tgpagella}
\usepackage{tgheros}
\usepackage[portuguese]{babel}
\usepackage{microtype}

% ---------- página e cores ----------
\usepackage[a4paper,top=2.6cm,bottom=2.4cm,left=2.4cm,right=2.4cm]{geometry}
\usepackage[table]{xcolor}
\definecolor{azul}{RGB}{28,62,112}
\definecolor{azulclaro}{RGB}{226,234,246}
\definecolor{cinza}{RGB}{95,102,115}

% ---------- tabelas, listas, caixas ----------
\usepackage{booktabs}
\usepackage{array}
\usepackage{tabularx}
\usepackage{enumitem}
\usepackage{titlesec}
\usepackage{fancyhdr}
\usepackage{lastpage}
\usepackage{tcolorbox}

\setlength{\parindent}{0pt}
\setlength{\parskip}{6pt}

\setlist{nosep, leftmargin=*, itemsep=3pt}
\setlist[itemize]{label=\textcolor{azul}{\textbullet}}
\setlist[enumerate]{label=\textcolor{azul}{\sffamily\bfseries\arabic*.}}

\titleformat{\section}{\sffamily\Large\bfseries\color{azul}}{\thesection}{0.7em}{}[{\color{azul}\titlerule[0.6pt]}]
\titleformat{\subsection}{\sffamily\large\bfseries\color{azul}}{\thesubsection}{0.6em}{}
\titlespacing*{\section}{0pt}{1.6em}{0.8em}
\titlespacing*{\subsection}{0pt}{1.1em}{0.3em}

\pagestyle{fancy}
\setlength{\headheight}{15pt}
\fancyhf{}
\renewcommand{\headrule}{{\color{azul}\hrule width\headwidth height\headrulewidth}}
\renewcommand{\headrulewidth}{0.4pt}
\fancyhead[L]{\sffamily\small\color{cinza} Universidade da Beira Interior}
\fancyhead[R]{\sffamily\small\color{cinza} Engenharia de Software \textperiodcentered{} 2026/27}
\fancyfoot[L]{\sffamily\small\color{cinza} Registo de Uso de IA}
\fancyfoot[R]{\sffamily\small\color{cinza} Página \thepage\ de \pageref{LastPage}}

\fancypagestyle{primeira}{%
  \fancyhf{}%
  \renewcommand{\headrulewidth}{0pt}%
  \fancyfoot[L]{\sffamily\small\color{cinza} Registo de Uso de IA}%
  \fancyfoot[R]{\sffamily\small\color{cinza} Página \thepage\ de \pageref{LastPage}}%
}

\newcolumntype{Y}{>{\raggedright\arraybackslash}X}
\newcolumntype{L}[1]{>{\raggedright\arraybackslash}p{#1}}
\newcolumntype{N}[1]{>{\raggedright\arraybackslash\sffamily\bfseries}p{#1}}
\newcommand{\th}[1]{\textcolor{white}{\sffamily\bfseries #1}}

\tcbset{
  caixa/.style={colback=white, colframe=azul, coltitle=white, colbacktitle=azul,
    fonttitle=\sffamily\bfseries, arc=3pt, boxrule=0.6pt,
    left=9pt, right=9pt, top=6pt, bottom=6pt}
}

% ---------- entradas do registo ----------
\newcommand{\linha}[1]{\rule[-2pt]{#1}{0.4pt}}
\newcommand{\claude}{Claude (Anthropic), via Claude Code}
% Se o campo vier vazio desenha uma linha para escrever à mão; se tiver texto, mostra o texto.
\newcommand{\campo}[2]{\def\conteudo{#1}\ifx\conteudo\empty\linha{#2}\else #1\fi}
% \entrada{data}{onde}{ferramenta}{para quê}{revisto por}{data da revisão}{}{}
% Os dois últimos ficam {}{} até alguém rever. Depois escreve-se lá o nome e a data.
\newcommand{\entrada}[6]{%
  \begin{tcolorbox}[caixa, title={#1 \textperiodcentered{} #2}]
    \renewcommand{\arraystretch}{1.45}%
    \begin{tabularx}{\linewidth}{@{}N{2.5cm} Y@{}}
      Ferramenta & #3 \\
      Para quê & #4 \\
      Revisto por & \campo{#5}{6.2cm}\hspace{1.2em}Data\hspace{0.5em}\campo{#6}{2.4cm} \\
    \end{tabularx}
  \end{tcolorbox}%
}
% Caixa em branco, para quem prefere preencher à mão depois de imprimir.
\newcommand{\entradavazia}{%
  \begin{tcolorbox}[caixa, title={Nova entrada}]
    \renewcommand{\arraystretch}{1.9}%
    \begin{tabularx}{\linewidth}{@{}N{2.5cm} Y@{}}
      Data & \linha{4cm} \\
      Onde & \linha{10.6cm} \\
      Ferramenta & \linha{10.6cm} \\
      Para quê & \linha{10.6cm} \\
      & \linha{10.6cm} \\
      Revisto por & \linha{6.2cm}\hspace{1.2em}Data\hspace{0.5em}\linha{2.4cm} \\
    \end{tabularx}
  \end{tcolorbox}%
}

\begin{document}
\thispagestyle{primeira}

% ---------- cabeçalho do documento ----------
\begin{tcolorbox}[colback=azul, colframe=azul, coltext=white, arc=3pt, boxrule=0pt,
    left=16pt, right=16pt, top=14pt, bottom=14pt]
  \sffamily
  {\footnotesize\bfseries\MakeUppercase{Universidade da Beira Interior}}\\[2pt]
  {\footnotesize Faculdade de Engenharia \textperiodcentered{} Departamento de Informática}\\[16pt]
  {\fontsize{28}{32}\selectfont\bfseries Registo de Uso de IA}\\[4pt]
  {\large Onde e para quê usamos IA no projeto \textperiodcentered{} Engenharia de Software 2026/27}
\end{tcolorbox}

Este ficheiro cumpre o que combinámos no Contrato de Equipa: dizer, de forma simples, onde e para quê usamos ferramentas de IA no projeto. Vamos atualizá-lo sempre que usarmos IA em algo novo. Está ligado no quadro Jira, na secção \textit{Artifacts}.

\begin{tcolorbox}[caixa, title={Declaração}, colback=azulclaro!60]
  A base do Contrato de Equipa foi feita com apoio de IA. Usámos o Claude, da Anthropic, através do Claude Code, para propor a estrutura do documento, escrever a primeira versão do texto em LaTeX, desenhar o aspeto do documento e sugerir alguns valores por omissão: a duração das reuniões, a rotação de Product Owner e Scrum Master e a escala de consequências.

  A partir daí, o contrato é da equipa. Lemos, ajustamos o que não fizer sentido para nós e respondemos por tudo o que lá está.
\end{tcolorbox}

%--------------------------------------------------------------------
\section{Registo}

\begin{tcolorbox}[caixa, title={Como se preenche}, colback=azulclaro!60]
  \begin{enumerate}
    \item Quem usou a IA acrescenta uma entrada nova no fim, com a data, o sítio onde usou, a ferramenta e para quê. No código há um modelo comentado: tira o \texttt{\%} do início das linhas e muda os textos. Há também caixas em branco no fim desta secção, para preencher à mão.
    \item Outro elemento da equipa lê o que a IA fez e confirma que está certo. Depois escreve o nome e a data nas duas últimas chavetas da entrada, no código.
    \item Enquanto essas chavetas estiverem vazias, aparece uma linha em branco e a entrada ainda não foi revista.
  \end{enumerate}
\end{tcolorbox}

% =====================================================================
% COMO ACRESCENTAR UMA ENTRADA
% 1. Copia as 7 linhas do modelo abaixo para o fim das entradas (antes das
%    caixas em branco) e tira o % do início de cada uma.
% 2. Muda só o que está dentro das chavetas. Não apagues as chavetas.
% 3. Deixa as duas últimas chavetas vazias: {}
% QUANDO ALGUÉM REVER: escreve o nome e a data nas duas últimas chavetas,
%    por exemplo {Nome Apelido} e {dd/mm/aaaa}.
% =====================================================================
%
% \entrada%
%   {dd/mm/aaaa} % data
%   {Onde} % onde (documento, código, Jira...)
%   {\claude} % ferramenta (\claude já escreve "Claude (Anthropic), via Claude Code")
%   {Para quê} % para quê
%   {} % revisto por
%   {} % data da revisão
%
% =====================================================================

\entrada%
  {09/10/2026} % data
  {Contrato de Equipa} % onde (documento, código, Jira...)
  {\claude} % ferramenta
  {Estrutura, primeira versão do texto, aspeto do documento e valores sugeridos (reuniões, rotação de PO e SM, consequências).} % para quê
  {Tiago Pereira} % revisto por (fica {} até alguém rever)
  {09/10/2026} % data da revisão (fica {} até alguém rever)

\entrada%
  {09/10/2026} % data
  {Registo de Uso de IA (este ficheiro)} % onde (documento, código, Jira...)
  {\claude} % ferramenta
  {Primeira versão do texto e, mais tarde, uma organização mais clara das entradas.} % para quê
  {} % revisto por (fica {} até alguém rever)
  {} % data da revisão (fica {} até alguém rever)

\entrada%
  {09/10/2026} % data
  {Guia de Jira} % onde (documento, código, Jira...)
  {\claude} % ferramenta
  {Pesquisa dos passos do Jira, primeira versão do guia e reescrita para a equipa começar a usar (também em página do Confluence).} % para quê
  {} % revisto por (fica {} até alguém rever)
  {} % data da revisão (fica {} até alguém rever)

\entrada%
  {09/10/2026} % data
  {Jira (estrutura e textos)} % onde (documento, código, Jira...)
  {\claude} % ferramenta
  {Primeira versão da estrutura e dos textos: epics, cartões, sprints, presenças e faltas, papéis, páginas do Confluence e ligação ao GitHub.} % para quê
  {} % revisto por (fica {} até alguém rever)
  {} % data da revisão (fica {} até alguém rever)

\entrada%
  {14/10/2026} % data
  {README e organização do repositório} % onde
  {Grok (xAI)} % ferramenta
  {Reescrita do README para se perceber o projeto, a equipa, os sprints e onde está cada documento. A pasta da análise do Sprint 1 passou a chamar-se sprint-1. No contrato, números de aluno e a parte do Git, que ainda falava na branch dev.} % para quê
  {} % revisto por (fica {} até alguém rever)
  {} % data da revisão (fica {} até alguém rever)

\entrada%
  {14/10/2026} % data
  {Análise do produto, Sprint 1} % onde
  {Grok (xAI)} % ferramenta
  {Leitura do enunciado (SE\_Project.pdf) e das aulas, e reescrita de docs/sprint-1/analise\_produto.tex. O produto passa a ser o sistema de pedidos de funcionalidades e roadmap. O Smart Campus do Lab03 fica como exercício. Saíram as personas, os cenários, as user stories e as prioridades da aplicação universitária. Essas secções ficam por preencher, com o formato das aulas. Fica a nota de que a apresentação do Sprint 1 e a do Sprint 2 são a mesma.} % para quê
  {} % revisto por (fica {} até alguém rever)
  {} % data da revisão (fica {} até alguém rever)

\entrada%
  {17/10/2026} % data
  {Análise do produto, Sprint 1} % onde
  {\claude} % ferramenta
  {Releitura do enunciado (SE\_Project.pdf) e primeira proposta de personas, cenários, user stories com critérios de aceitação, requisitos não funcionais e interfaces candidatas, em docs/sprint-1/analise\_produto.tex. As personas são inventadas e as prioridades (MoSCoW) são só uma proposta. O levantamento com o Character.ai ainda não foi feito.} % para quê
  {} % revisto por (fica {} até alguém rever)
  {} % data da revisão (fica {} até alguém rever)

\entrada%
  {17/10/2026} % data
  {Documentos dos Sprints 1 e 2 em PDF} % onde
  {\claude} % ferramenta
  {Paginação em PDF da análise do produto (docs/sprint-1/produto\_personas\_cenarios.pdf), do guião das entrevistas no Character.ai com as folhas de respostas (docs/sprint-2/entrevistas\_character\_ai.pdf) e dos rascunhos dos ecrãs (docs/sprint-2/wireframes\_exemplos.pdf).} % para quê
  {} % revisto por (fica {} até alguém rever)
  {} % data da revisão (fica {} até alguém rever)

\entrada%
  {17/10/2026} % data
  {Wireframes no Figma} % onde
  {\claude} % ferramenta
  {Quinze ecrãs no ficheiro Sprint 2 do Figma: uma versão em estilo de quadro Scrum (Quadro v2) e a versão de baixa fidelidade, só com caixas e sem texto (Quadro v3), a partir dos rascunhos e das áreas do enunciado.} % para quê
  {} % revisto por (fica {} até alguém rever)
  {} % data da revisão (fica {} até alguém rever)

\entrada%
  {17/10/2026} % data
  {Jira, personas, cenários e user stories} % onde
  {\claude} % ferramenta
  {Passar para o Jira o que saiu das entrevistas que a equipa fez no Character.ai (capturas e texto no cartão EDS-55), ligar personas a cenários e cenários a user stories, e dar nomes em português aos epics. As entrevistas foram feitas pela equipa, não pela IA.} % para quê
  {} % revisto por (fica {} até alguém rever)
  {} % data da revisão (fica {} até alguém rever)

\entrada%
  {19/10/2026} % data
  {Jira, cartões dos wireframes} % onde
  {\claude} % ferramenta
  {Em cada cartão dos wireframes (EDS-97 a EDS-103), uma tabela com os problemas das personas que o ecrã resolve e o que ainda falta no Figma, a partir das entrevistas.} % para quê
  {} % revisto por (fica {} até alguém rever)
  {} % data da revisão (fica {} até alguém rever)

\entrada%
  {23/10/2026} % data
  {Jira, Sprint 3, e README} % onde
  {\claude} % ferramenta
  {Proposta do sprint backlog do Sprint 3 e do Sprint 4 (cartões, critérios de aceitação, estimativas), proposta de arquitetura (EDS-153), cartão da apresentação oral (EDS-173) e README atualizado. As decisões ficam para a planning com o Product Owner.} % para quê
  {} % revisto por (fica {} até alguém rever)
  {} % data da revisão (fica {} até alguém rever)

\entrada%
  {23/10/2026} % data
  {Wireframes, diagramas, código e wiki do Sprint 3} % onde
  {\claude} % ferramenta
  {Correção de 7 ecrãs dos wireframes de baixa fidelidade com o que faltava das personas; diagramas no FigJam (arquitetura, casos de uso, duas sequências, estados, atividades, modelo de dados); base do código (Django REST Framework e React, 15 ecrãs sem cores) com 28 testes no back-end e 23 no front-end; integração contínua; AGENTS.md e a pasta wiki/ em Markdown; exportação dos wireframes, fontes Mermaid dos diagramas e documento de requisitos não funcionais em LaTeX (docs/sprint-3/). As ideias foram da equipa, na reunião de 23/10.} % para quê
  {} % revisto por (fica {} até alguém rever)
  {} % data da revisão (fica {} até alguém rever)

\entrada%
  {24/10/2026} % data
  {Estilo da aplicação e Docker} % onde
  {\claude} % ferramenta
  {CSS simples da aplicação (cores neutras, menu com o ecrã atual, filtros numa linha); Dockerfile do back-end e do front-end e docker-compose.yml com os dois contentores na mesma rede; correção da proteção CSRF para os pedidos vindos do front-end, com teste novo.} % para quê
  {} % revisto por (fica {} até alguém rever)
  {} % data da revisão (fica {} até alguém rever)

\entrada%
  {24/10/2026} % data
  {README e contas de teste} % onde
  {\claude} % ferramenta
  {README reescrito com prints da aplicação e dos diagramas (tirados da aplicação a correr e dos ficheiros .mmd de docs/sprint-3/modelacao); palavras-passe das contas de teste passaram a ser o nome seguido de 123, com uma conta admin.} % para quê
  {} % revisto por (fica {} até alguém rever)
  {} % data da revisão (fica {} até alguém rever)

\entradavazia

\entradavazia

%--------------------------------------------------------------------
\section{Ferramentas}

\small
\renewcommand{\arraystretch}{1.5}
\begin{tabularx}{\textwidth}{N{4.4cm} Y L{3.2cm}}
  \rowcolor{azul} \th{Ferramenta} & \th{Para que a usamos} & \th{Estado} \\
  \rowcolor{azulclaro!55}
  Claude (Anthropic), via Claude Code & Texto e LaTeX dos documentos da equipa. & Em uso \\
  \rowcolor{white}
  Grok (xAI) & README e acertos pontuais nos documentos da equipa. & Em uso \\
  \rowcolor{azulclaro!55}
  Character.ai & Levantar requisitos, como o enunciado pede. & Em uso: entrevistas às cinco personas a 17/10 (EDS-55) \\
  \rowcolor{white}
  Ferramenta de IA para programar & Apoio ao código. A equipa escolhe pelo menos uma e regista aqui qual é e porquê. & Por escolher \\
  \bottomrule
\end{tabularx}
\renewcommand{\arraystretch}{1}
\normalsize

%--------------------------------------------------------------------
\section{O que registamos}

\begin{itemize}
  \item Tudo o que a IA tenha influenciado naquilo que entregamos: texto, código, diagramas, requisitos, personas, testes.
  \item Em cada linha: a data, onde foi usada, a ferramenta, para quê e quem reviu o resultado.
  \item Quando for código, juntamos a ligação para o \textit{pull request}.
  \item Dúvidas pontuais e pesquisas que não mudaram nada no que entregamos não precisam de linha.
\end{itemize}

\section{Como garantimos que a IA não decide por nós}

\begin{itemize}
  \item Só entra no projeto o que percebemos e conseguimos explicar nas apresentações. Cada um responde pelo que entrega.
  \item O que é gerado por IA passa pela mesma revisão, pelos mesmos testes e pela mesma Definição de Concluído que o resto.
  \item Não pomos passwords, chaves, dados pessoais nem informação que não possamos partilhar em ferramentas de IA.
  \item Verificamos a informação, as fontes e as licenças do que a IA nos sugere.
\end{itemize}

\end{document}
